\chapter{Rendimiento y Optimización del Puente}
\label{chap:performance}
{\huge$\bigstar$}{\huge$\bigstar$} Intermedio

\section{Overhead de py:call}
\begin{enumerate}
\item GIL acquire/release ($\sim\mu$s)
\item Conversión Lisp $\to$ Python
\item Conversión Python $\to$ Lisp
\item Refcounting
\end{enumerate}

\section{Regla Práctica}
\textbf{$<$ 1 000 llamadas:} insignificante.

\textbf{$>$ 10 000 llamadas:} usar \texttt{py:run} por lotes.

\section{Modos de Uso}
\paragraph{Modo 1: Lisp Nativo + Python Puntual}
Control en Lisp, Python para operaciones específicas.

\paragraph{Modo 2: Python Delegado}
\texttt{py:run} para bloques intensivos.

\paragraph{Modo 3: VALUE\_PYTHON Directo}
Llamada directa a callable, sin \texttt{py:call}.

\begin{examplebox}[title=Benchmark comparativo]
\begin{lispcode}
(defun test-pycall ()
  (let ((s (py:attr (py:import "math") "sqrt")))
    (time (dotimes (i 10000) (s (float i))))))
\end{lispcode}
\end{examplebox}

